% ============================================================
%  Kapitel 5: Implementierung und Entwicklung
%  ARis -- AR-Brille Studienarbeit | Bastian Kiefer
%  Überarbeitete Fassung, April 2026
% ============================================================

\chapter{Implementierung und Entwicklung}
\label{chap:implementierung}

% ---------------------------------------------------------------------------
\section{Optische Implementierung}
\label{sec:optischeImplementierung}

Der in Abschnitt~\ref{subsec:FreeSpaceCombiner} beschriebene Free-Space-Combiner-Aufbau
wurde mit bewusst kostengünstigen Komponenten realisiert, um die Machbarkeit des
Gesamtsystems innerhalb der gegebenen Budgetgrenzen nachzuweisen.
Als Lösung für das Mikrodisplay-Linsensystem erwies sich der elektronische
Sucher (\ac{EVF}) einer modernen Digitalkamera als geeignete, sofort einsatzbereite
Einheit.
Viewfinder moderner Systemkameras integrieren ein Micro-OLED-Display sowie ein
abgestimmtes Okular in einem kompakten Gehäuse und sind damit funktional äquivalent
zu eigens konfigurierten NED-Komponenten, jedoch zu einem Bruchteil der Kosten.

Zu Beginn der Projektlaufzeit wurde die Beschaffung einer Waveguide-Linse angestrebt,
welche aufgrund ihrer Dünnheit den Innovationscharakter des Systems erheblich
gesteigert hätte.
Entsprechende Anfragen bei Zulieferern ergaben jedoch Mindestpreise von
ca.\ 1.000,-€ pro Einheit, sodass diese Option im Rahmen des Projektbudgets
verworfen wurde.
Der halbdurchlässige Combiner-Spiegel wurde als handelsüblicher Einwegspiegel
beschafft, der die in Abschnitt~\ref{subsec:FreeSpaceCombiner} beschriebene
See-Through-Eigenschaft des optischen Systems bei geringen Kosten hinreichend
realisiert.
Nach Eingang aller Komponenten wurden diese in den mechanischen Trägerrahmen
integriert und die optische Funktion des Gesamtaufbaus qualitativ und iterativ
verifiziert.

Ein unerwartetes Ergebnis dieser Verifikationsphase war die geringe effektive
Vergrößerung des gewählten Viewfinder-Okulars: Die tatsächlich wahrgenommene
Bildgröße auf dem Combiner-Spiegel fiel deutlich kleiner aus als erwartet worden war, sodass das Bild lediglich
einen kleinen Ausschnitt des Sichtfeldes bedeckte.
Dies ist auf die kurze Brennweite des verbauten Okulars zurückzuführen,
die zwar eine entspannte Augenakkommodation ermöglicht, jedoch das
erreichbare Field of View im gewählten Aufbau begrenzt (vgl.~Abschnitt~\ref{subsec:FoV}).
\newpage
Mit mehr Zeit und Budget könnte hier durch ein Okular mit angepasster Brennweite das tatsächlich angezeigte Bild
deutlich vergrößert und damit der konzipierte Anzeigebereich vollständig genutzt werden. Ebenso würde ein längerer Strahlengang zu einem größeren Bild führen. Dies war Anfangs von der Optik auch vorgesehen (wie in \ref{img:optischesgesamtkonzept} gezeigt), wurde jedoch im mechanischen Design nicht so umgesetzt. Das seitliche Design wurde durch einen verkleinerten frontalen Aufsatz ersetzt, was der mobilität hilft, jedoch die Ausbreitung des Bildes verhindert.

Als pragmatische Konsequenz dieser Erkenntnis wurde eine dedizierte Demo-Seite
in der Onboard-UI entwickelt (vgl.~Abschnitt~\ref{subsec:demoSeite}), die den
verfügbaren Anzeigebereich durch eine kompakte, informationsdichte Darstellung
optimal ausnutzt.
Das vollständige Benutzerinterface lässt sich alternativ über einen HDMI-Ausgang
an einem externen Monitor darstellen und testen, da der Raspberry~Pi~Zero~2\,W
einen entsprechenden Ausgang bereitstellt.

% ---------------------------------------------------------------------------
\section{Software-Implementierung}
\label{sec:softwareImplementierung}

Die Softwareentwicklung verlief nicht als linearer Prozess, sondern reflektiert
die realen Randbedingungen eines iterativen Prototypenprojekts.
In einer ersten Phase wurden Frontend-Grundstruktur und Backend-Konzept weitgehend
parallel und zunächst systemgetrennt -- ohne physische Hardware -- entwickelt.
Dabei wurden auch verschiedene Technologiealternativen evaluiert: Ein früher
Frontend-Prototyp entstand als Vue.js-Webapp
(\texttt{Companion\_Prototype/Vue\_Companion\_App/}), die grundlegende
REST-Kommunikation mit dem Backend erprobte, bevor der endgültige
SvelteKit-Stack festgelegt wurde.

Mit Eingang des Raspberry~Pi und des Displays begann die zweite Phase, in der
Hardware-Tests durchgeführt und das Backend auf die reale Zielplattform abgestimmt
wurden.
Die Hardwarekomponenten wurden anschließend an die für die Mechanik verantwortliche
Projektpartnerin übergeben, was die weitere Softwareentwicklung temporär unterbrach.
In der abschließenden Integrationsphase -- in der Woche vor der Abgabe -- wurde
das Gesamtsystem nach Rückgabe der Brille zusammengeführt und auf dem Raspberry~Pi
deployed.

Diese letzte Phase brachte unerwartete Deployment-Herausforderungen mit sich:
Neuere Builds des Raspberry~Pi~OS deaktivieren standardmäßig den analogen
AV-Ausgang, was die Ansteuerung des verbauten Displays verhinderte.
Ein Downgrade auf ein älteres Image war notwendig, das wiederum ältere
Paketversionen -- darunter Python~3.9 und ältere \texttt{npm}-Versionen --
mitbrachte, welche Kompatibilitätsprobleme mit den verwendeten Bibliotheken
verursachten.
Die Auflösung dieser Abhängigkeitskonflikte erforderte erheblichen Aufwand;
die resultierende, reproduzierbare Installationsanleitung ist im begleitenden
GitHub-Repository dokumentiert.

Der knappe Zeitrahmen der Integrationsphase führte dazu, dass der Fokus auf
die Vervollständigung des Grundsystems -- stabiler Start, funktionierender
Datenpfad vom Backend zum Display -- gelegt wurde, anstatt alle konzipierten
Features vollständig zu finalisieren.
Auf die konkreten Auswirkungen dieser Entscheidungen auf den Funktionsumfang
wird in Abschnitt~\ref{chap:Diskussion} eingegangen.

Der gesamte Code wurde unter Einsatz moderner KI-gestützter Entwicklungswerkzeuge
erstellt.
Tools wie ChatGPT Codex und vergleichbare Assistenzsysteme ermöglichen es,
auch in weniger vertrauten Frameworks produktiv zu arbeiten, indem sie
Boilerplate-Code generieren, Bibliotheks-APIs erklären und beim Debugging
unterstützen.
Dies ist als genuiner Vorteil moderner Softwareentwicklung zu verstehen:
Die eingesetzte Technologiekombination -- Svelte, Flask, React Native -- hätte ohne diese Unterstützung selbst für einen gelernten Softwareentwickler eine Herausforderung dargestellt und somit konnte der Fokus weniger auf Wiederholungen von Basis-Code und mehr auf die spezifischen Herausforderungen des Gesamtsystems fokussiert werden.
Der vollständig dokumentierte Quellcode ist im begleitenden GitHub-Repository
einsehbar; im vorliegenden Kapitel wird der Code ausschließlich prinzipiell
und exemplarisch behandelt.
\newpage
% ---------------------------------------------------------------------------
\subsection{Modulübersicht}
\label{subsec:moduluebersicht}

Das Gesamtsystem besteht aus vier klar getrennten Softwaremodulen
(vgl.\ Abbildung~\ref{fig:modulmap}):

\begin{enumerate}
  \item \textbf{Backend} (\texttt{Software/Backend/}) -- Flask-basierter
    Python-Server auf dem Raspberry~Pi, der alle Dienste koordiniert und die
    REST-API bereitstellt.
  \item \textbf{Onboard-UI} (\texttt{Software/Onboard\_UI/}) -- SvelteKit-Applikation,
    die als lokaler Webserver auf dem Raspberry~Pi läuft und die AR-Anzeige
    auf dem Display ausgibt.
  \item \textbf{Onboard-Runtime} (\texttt{Software/Onboard\_Runtime/}) -- Python-Paket,
    das gemeinsam genutzte Laufzeitdienste kapselt: Modulverwaltung, Ereignisbus,
    Kommunikationsadapter und Hardware-Abstraktion.
  \item \textbf{Companion-App} (\texttt{Software/Aris\_Companion/}) -- React-Native-
    Anwendung auf dem Mobilgerät des Nutzers zur Konfiguration und Fernsteuerung.
\end{enumerate}

\newpage
\begin{landscape}
\begin{figure}[htbp]
  \centering
  \begin{tikzpicture}[
      node distance=0.6cm and 1.4cm,
      box/.style={draw, rounded corners=3pt, minimum width=3.0cm,
                  minimum height=0.65cm, font=\small, fill=green!8},
      grp/.style={draw, dashed, rounded corners=4pt, inner sep=6pt,
                  fill=gray!5},
      arr/.style={->, >=stealth, thick},
      lbl/.style={font=\footnotesize\itshape}
    ]
    % --- Backend (links) ---
    \node[box] (flask)    {Flask App};
    \node[box, below=of flask]   (router)   {Route-Blueprints};
    \node[box, below=of router]  (services) {Services};
    \node[box, below=of services](repo)     {SQLiteRepository};
    \node[box, below=of repo]    (sqlite)   {SQLite (data.db)};
    \draw[arr] (flask)--(router);
    \draw[arr] (router)--(services);
    \draw[arr] (services)--(repo);
    \draw[arr] (repo)--(sqlite);

    % --- Onboard-Runtime (Mitte) ---
    \node[box, right=1.8cm of flask]  (runtime) {OnboardRuntime};
    \node[box, below=of runtime]      (ebus)    {SharedEventBus};
    \node[box, below=of ebus]         (comm)    {Comm-Adapter (BT / WiFi / USB)};
    \node[box, below=of comm]         (hw)      {Hardware-Adapter (I2C / GPIO)};
    \draw[arr] (runtime)--(ebus);
    \draw[arr] (ebus)--(comm);
    \draw[arr] (ebus)--(hw);

    % --- Onboard-UI (rechts) ---
    \node[box, right=1.8cm of runtime] (svelte) {SvelteKit (Onboard-UI)};
    \node[box, below=of svelte]        (fhost)  {FeatureHost};
    \node[box, below=of fhost]         (input)  {InputController (Taster-FSM)};
    \draw[arr] (svelte)--(fhost);
    \draw[arr] (fhost)--(input);

    % --- Companion (ganz rechts) ---
    \node[box, right=1.8cm of svelte] (companion) {Companion App (React Native)};

    % --- Querverbindungen ---
    \draw[arr, dashed] (svelte.west) -- node[above,lbl]{REST} (runtime.east);
    \draw[arr, dashed] (runtime.west) -- node[above,lbl]{REST} (flask.east);
    \draw[arr, dashed] (companion.west) -- node[above,lbl]{REST / BLE}
      (svelte.east);
  \end{tikzpicture}
  \caption{Modulkarte des ARis-Softwaresystems.}
  \label{fig:modulmap}
\end{figure}
\end{landscape}

Die Kommunikationspfade folgen dem in Kapitel~\ref{sec:Software-Architektur} beschriebenen
Architekturprinzip der losen Kopplung: Jedes Modul kommuniziert ausschließlich
über definierte Schnittstellen mit den anderen, was eine unabhängige Entwicklung
und Erweiterung der einzelnen Komponenten ermöglicht.

% ---------------------------------------------------------------------------
\subsection{Onboard-Runtime}
\label{subsec:onboardRuntime}

Die \texttt{Onboard\_Runtime} ist ein eigenständiges Python-Paket, das
systemübergreifende Laufzeitdienste kapselt, die weder zum Backend noch zur UI
gehören.
Sie bildet das verbindende Rückgrat, auf das alle anderen Onboard-Module
zurückgreifen können, ohne direkte Abhängigkeiten untereinander einzugehen.

Kernkomponente ist der \texttt{SharedEventBus} (\texttt{event\_bus.py}),
ein In-Prozess-Nachrichtenbus nach dem Publish/Subscribe-Muster.
Dabei abonnieren Module Ereignisse eines bestimmten Typs -- etwa
\texttt{BUTTON\_PRESS} oder \texttt{SENSOR\_UPDATE} -- und werden bei deren
Auftreten aufgerufen, ohne den Sender kennen zu müssen.
Dieses Entkopplungsprinzip trennt die Hardwareschicht vollständig von der
Anzeigelogik: Ein Tasterdruck löst ein Ereignis aus, das der EventBus an
alle registrierten Empfänger weiterleitet, ohne dass der Taster-Code
wissen muss, welche Funktion er damit auslöst.

Die \texttt{OnboardRuntime}-Klasse (\texttt{runtime.py}) orchestriert den
Startablauf und verwaltet alle Kernmodule über einen einheitlichen Lebenszyklus
mit den Phasen \texttt{init}, \texttt{start}, \texttt{stop} und \texttt{health}.
Startup-Profile (\texttt{dev}, \texttt{demo}, \texttt{headless}) ermöglichen
einen kontextabhängigen Betrieb, beispielsweise für Entwicklungsumgebungen ohne
angeschlossene Hardware.

Das Kommunikationssubsystem (\texttt{comm/}) enthält Adapter für Bluetooth
(GATT-Server), WiFi sowie einen USB-Debug-Adapter für die Entwicklungsphase.
Ein \texttt{TransportTranslator} normalisiert die unterschiedlichen
Transportformate auf ein einheitliches internes Kommandoformat, sodass höhere
Schichten transportagnostisch implementiert werden können und ein späterer
Wechsel des Kommunikationskanals -- etwa von WiFi zu Bluetooth -- keine
Änderungen in der Anwendungslogik erfordert.

% ---------------------------------------------------------------------------
\subsection{Python-Backend}
\label{subsec:backendImpl}

Das Python-Backend ist als Flask-Applikation realisiert und bildet die zentrale
Koordinationsschicht des Systems.
Um eine klare Trennung der Verantwortlichkeiten zu erreichen, folgt es einer
dreischichtigen Struktur: Route-Handler nehmen HTTP-Anfragen entgegen und
delegieren an Service-Klassen, die ihrerseits ausschließlich über eine
\texttt{SQLiteRepository}-Klasse auf die Datenbank zugreifen.
Jede dieser Schichten kann damit unabhängig getestet und ausgetauscht werden,
ohne die anderen Schichten zu berühren.

Um Seitenstartcode zu vermeiden und das spätere Testen einzelner Komponenten zu
vereinfachen, wird die Flask-Anwendung über eine Factory-Funktion
(\texttt{create\_app()} in \texttt{app.py}) erzeugt, die alle Abhängigkeiten
konfiguriert und die Route-Blueprints registriert.
Blueprints sind Flask-eigene Module zur Gruppierung thematisch zusammengehöriger
Endpunkte -- beispielsweise alle Teleprompter-bezogenen Routen in einem
\texttt{teleprompter\_bp}-Blueprint -- und ermöglichen so eine übersichtliche
Strukturierung des Backend-Codes.

Tabelle~\ref{tab:endpoints} gibt einen Überblick über die implementierten
REST-API-Endpunkte des Systems.
Die Endpunkte decken alle konzipierten Funktionsbereiche ab: Systemstatus und
Sensorik, persistente Nachrichten, Teleprompter-Steuerung, Navigation sowie
Geräteeinstellungen.

\begin{table}[htbp]
  \centering
  \caption{Implementierte REST-API-Endpunkte des ARis-Backends.}
  \label{tab:endpoints}
  \begin{tabularx}{\textwidth}{@{} l l X @{}}
    \toprule
    \textbf{Endpunkt} & \textbf{Methode(n)} & \textbf{Funktion} \\
    \midrule
    \texttt{/api/mainInfo}               & GET              & Akkustand, Temperatur, Luftfeuchtigkeit \\
    \texttt{/api/status}                 & GET              & Systemstatus, Hardware-Readiness \\
    \texttt{/api/sensors}                & GET              & I2C-Sensor-Snapshot (MPU-6050, GY-63) \\
    \texttt{/api/messages}               & GET, POST, DELETE & Persistente Nachrichten \\
    \texttt{/api/teleprompter/send}      & POST             & Neue Konfiguration von Companion App \\
    \texttt{/api/teleprompter/current}   & GET              & Polling-Endpunkt für AR-Brille \\
    \texttt{/api/teleprompter/history}   & GET              & Letzte gespeicherte Konfigurationen \\
    \texttt{/api/teleprompter/reset}     & POST             & Zurücksetzen auf Standardwerte \\
    \texttt{/api/navigation/current}     & GET              & Orientierungsdaten aus MPU-6050 \\
    \texttt{/api/settings/device}        & GET, PUT         & Geräteeinstellungen \\
    \texttt{/health}                     & GET              & Liveness-Probe \\
    \bottomrule
  \end{tabularx}
\end{table}

\subsubsection{Hardware-Abstraktion und Sensorik}
\label{subsubsec:hardwareImpl}

Die Hardwareanbindung ist vollständig im \texttt{SensorService} gekapselt,
der als einzige Schicht direkten Zugriff auf die GPIO- und I2C-Schnittstellen
des Raspberry~Pi hat.
Der Dienst abstrahiert drei Komponenten: den physischen Taster
(Board-Pin~40) mit Software-Entprellung, den MPU-6050 über I2C für
Beschleunigung und Gyroskop sowie den GY-63/BMP085 über I2C als
Umgebungssensor für Temperatur und Luftfeuchtigkeit.
Alle Hardware-Pin-Zuweisungen sind in einer zentralen \texttt{pinmap.py} definiert,
sodass eine Portierung auf abweichende Hardware keine Änderungen im
Service-Code erfordert -- lediglich die Pinbelegung muss angepasst werden.

Besonders praktisch in der Entwicklung: Die Erweiterbarkeit des Systems durch weitere Sensoren, wie zum Beispiel ein BME280, wird durch das modulare Design vereinfacht. Mit dem passenden Eintrag in die \texttt{pinmap.py} und die entsprechende Erweiterung von \texttt{sensor\_service.py} kann das Auslesen eines neuen Sensors und dessen Verarbeitung einfach ermöglicht werden. Die Änderungen des \texttt{SensorService} wird der Sensor REST-Endpoint erweitert und benötigt keine tiefgründigeren Änderungen. Anschließend können die Daten auf der gewünschten Seite abgefragt werden.

\subsection{Datenhaltung}
\label{subsec:datenbankImpl}

Die Datenhaltung erfolgt über eine eingebettete SQLite-Datenbank
(\texttt{data.db}), die ohne separaten Datenbankserver direkt im Dateisystem
des Raspberry~Pi persistiert wird.
SQLite wurde aufgrund seiner geringen Ressourcenanforderungen gewählt, die dem
eingeschränkten Rechenumfeld des Raspberry~Pi~Zero~2\,W entgegenkommen.
Der gesamte Datenbankzugriff erfolgt über die Klasse \texttt{SQLiteRepository},
die alle Operationen hinter semantischen Methoden kapselt und die
Anwendungslogik damit vollständig von der konkreten Speichertechnologie trennt.
Das Schema umfasst vier Tabellen mit klar abgegrenzten Aufgaben: 

\begin{itemize}
  \item \textbf{messages} speichert eingehende Textnachrichten, die über
    die Companion App übermittelt und auf der AR-Brille angezeigt werden.
    Die Tabelle kann über den DELETE-Endpunkt vollständig geleert werden.
  \item \textbf{teleprompter\_history} protokolliert jede von der Companion App
    gesendete Teleprompter-Konfiguration mit vollständigem Parametersatz
    (Text, Geschwindigkeit, Schriftgröße, Farbe, Schriftart, Zeilenabstand,
    Transparenz) und Zeitstempel, was die Wiederverwendung früherer Einstellungen
    ermöglicht.
  \item \textbf{pairing\_sessions} verwaltet laufende und abgeschlossene
    Pairing-Vorgänge zwischen Companion App und Backend.
  \item \textbf{auth\_tokens} speichert ausgestellte Authentifizierungstoken
    als Hash und ermöglicht eine gesicherte Kommunikation
    zwischen Companion und Backend.
\end{itemize}

% ---------------------------------------------------------------------------
\subsection{Frontend -- SvelteKit Onboard-UI}
\label{subsec:frontendImpl}

Das AR-Frontend ist als SvelteKit-Applikation realisiert, die auf dem
Raspberry~Pi als lokal servierter Webserver läuft und deren Ausgabe direkt
auf dem im Brillengehäuse verbauten Display dargestellt wird.
Das visuelle Design orientiert sich an den Anforderungen eines
Head-up-Displays: Grüne Schrift (\texttt{\#0f0}) auf schwarzem Hintergrund
(\texttt{\#000}) stellt einen hohen Kontrast zum realen Umgebungsdurchblick
sicher und minimiert die Verwechslung virtueller Inhalte mit realen Szenen.
Alle Farbwerte sind über CSS-Variablen in einer globalen Stylesheet-Datei
zentralisiert und können damit systemweit konsistent angepasst werden.

\subsubsection{Aufbau einer Svelte-Komponente}
\label{subsubsec:svelteAufbau}

SvelteKit organisiert die Anwendung über ein Dateisystem-basiertes Routing:
Jeder Unterordner unter \texttt{src/routes/} entspricht einer eigenständigen
Route, die über den Pfad \texttt{/[Ordnername]} im Browser erreichbar ist.
Jede Seite besteht aus einer \texttt{+page.svelte}-Datei, die den typischen
dreiteiligen Aufbau einer Svelte-Komponente trägt (siehe \ref{lst:svelteHeader}).
Der \texttt{<script>}-Block enthält die TypeScript-Logik; der Markup-Teil darunter
definiert die Darstellung.
Reaktive Datenbindungen werden in Svelte durch einfache Variablenzuweisung
ausgedrückt -- ohne separates State-Management-Framework.
Sobald sich der Wert von \texttt{data} ändert, aktualisiert Svelte automatisch
nur die betroffenen DOM-Stellen.
Eine globale Layout-Datei (\texttt{+layout.svelte}) bindet den beschriebenen
Header systemweit ein, sodass er über alle Seiten hinweg konsistent erscheint.

\newpage
\begin{lstlisting}[language=HTML, caption={Grundstruktur einer Svelte-Komponente am Beispiel des System-Headers.}, label=lst:svelteHeader]
<script lang="ts">
  import { onMount } from 'svelte';

  // reaktive Zustandsvariable
  let data: { Battery?: string; Temperature?: string; Humidity?: string } = {};

  async function fetchData() {
    const res = await fetch('http://localhost:5000/api/mainInfo');
    data = await res.json();
  }

  // onMount: wird einmalig nach dem Einbinden der Komponente ausgefuehrt
  onMount(() => {
    fetchData();
    const interval = setInterval(fetchData, 5000);
    return () => clearInterval(interval); // Cleanup beim Entfernen
  });
</script>

<!-- Markup-Block: reaktiv gebunden an "data" -->
<div id="interface">
  <h1>
    Akku: {data.Battery ?? '?'} |
    Temp: {data.Temperature ?? '?'} |
    Hum: {data.Humidity ?? '?'}
  </h1>
</div>
\end{lstlisting}
\newpage

\subsubsection{Implementierte Seiten}
\label{subsubsec:implementierteSeiten}

Im finalen Stand sind die folgenden Seiten implementiert:

\begin{itemize}
  \item \textbf{Hub} (\texttt{/}) -- Startseite mit karussellartiger
    App-Auswahl; Einstiegspunkt für alle weiteren Seiten.
  \item \textbf{Teleprompter} (\texttt{/Teleprompter}) -- Scrollender Text
    mit vollständig konfigurierbaren Parametern (Geschwindigkeit, Schriftgröße,
    Farbe); die Konfiguration wird von der Companion App übermittelt und
    sekündlich vom Backend abgerufen.
    Der Teleprompter ist funktionsfähig implementiert; aufgrund des
    beschränkten Zeitrahmens in der Integrationsphase wurde jedoch kein
    vollständiger Systemtest auf der finalen Hardware durchgeführt.
  \item \textbf{Messages} (\texttt{/Messages}) -- Anzeige persistierter Nachrichten, welche
    mit der Companion App in einem KI-Chat erhalten wurden. Diese Seite ist vorbereitet, bietet jedoch noch keine implementierte Funktionalität.
  \item \textbf{Navigation} (\texttt{/Navigation}) -- Darstellung von
    Orientierungsdaten aus dem MPU-6050 sowie Navigationsdaten der Companion App; die Seite ist strukturell vorbereitet, eine vollständige Navigationsfunktion wurde im Rahmen
    des Projekts jedoch nicht implementiert.
  \item \textbf{Debug} (\texttt{/Debug}) -- Hardware- und Sensor-Diagnosepanel,
    das den vollständigen Systemstatus des Backends visualisiert und für den
    Produktionsbetrieb deaktiviert werden kann.
  \item \textbf{Demo} (\texttt{/Demo}) -- kompakt gestaltete Seite, die als
    Reaktion auf die unerwartet kleine Bildgröße des Viewfinder-Okulars
    entwickelt wurde (vgl.~Abschnitt~\ref{sec:optischeImplementierung}).
    Sie zeigt die wichtigsten Systemdaten -- Akkustand, Sensorwerte,
    Zeit -- auf minimalem Bildschirmraum und demonstriert gleichzeitig
    die Interaktionsmöglichkeit über den physischen Taster.
    \label{subsec:demoSeite}
\end{itemize}

Neue Seiten lassen sich durch Anlegen eines neuen Ordners mit einer
\texttt{+page.svelte}-Datei unter \texttt{src/routes/} hinzufügen, ohne
weitere Konfigurationsänderungen vornehmen zu müssen. Unter \ref{apx:implementierteseiten} finden sich Screenshots einiger implementierter Seiten.

\subsubsection{Eingabesteuerung}
\label{subsubsec:inputController}

Da das Brillengehäuse keinen Touchscreen besitzt, basiert die gesamte
Nutzerinteraktion auf einem einzigen physischen Taster.
Das Modul \texttt{input-controller.ts} implementiert hierfür eine
Finite-State-Machine, die verschiedene Eingabemuster -- Single-Tap und
Double-Tap -- kontextabhängig unterschiedlich interpretiert.
Im Hub-Kontext wechselt ein Single-Tap zwischen den verfügbaren Apps;
ein Double-Tap wählt die entsprechende App aus.
Im Feature-Kontext werden seitenspezifische Aktionen ausgelöst, etwa das
Starten oder Stoppen des Teleprompters.
Eingabeereignisse werden als Browser-Custom-Events über \texttt{window.dispatchEvent()}
propagiert, sodass jede Seite eigene Handler registrieren kann, ohne direkt
mit dem Input-Controller gekoppelt zu sein.
Für Entwicklungsumgebungen ohne physischen Taster ist eine
Tastatur-Emulation implementiert.

% ---------------------------------------------------------------------------
\subsection{Companion-App}
\label{subsec:companionImpl}

Die Companion App dient als mobile Steuerzentrale, über die der Nutzer alle
AR-Features konfigurieren und steuern kann, ohne direkt mit dem Raspberry~Pi
interagieren zu müssen.
Im bisherigen Entwicklungsstand kommuniziert die App primär über WiFi mit dem
Flask-Backend, da beide Geräte dabei im selben Netzwerk erreichbar sein müssen
und die REST-Schnittstelle so ohne zusätzliche Konfiguration genutzt werden kann.
Diese Lösung war für die Prototypenentwicklung funktional ausreichend, bringt
jedoch die Einschränkung mit sich, dass der Raspberry~Pi zunächst über eine
SSH-Verbindung oder eine lokale Session mit einem Netzwerk verbunden werden muss und keine einfache Input möglichkeit für die initiale Einrichtung zur Verfügung steht.
\newpage
Mittelfristig ist der Übergang zu einer primär Bluetooth-basierten Verbindung
vorgesehen: Das Backend verfügt bereits über einen implementierten
GATT-Server-Adapter in der Onboard-Runtime, der eine direkte BLE-Verbindung
zwischen Smartphone und Raspberry~Pi ermöglicht -- unabhängig von einer
gemeinsamen WLAN-Infrastruktur. Diese Bluetoothschnittstelle wurde bereits getestet, jedoch wegen des Zeitrahmens nicht final implementiert.
Die finale Companion App ist als native React-Native-Anwendung ausgeführt,
da nur ein nativer App-Kontext den erforderlichen Zugriff auf die
Bluetooth-Low-Energy-APIs des Betriebssystems erlaubt.
Im Laufe der Entwicklung wurde zunächst ein Vue.js-basierter Web-Prototyp
eingesetzt, der die REST-Kommunikation mit dem Backend erprobte, bevor der
endgültige Stack festgelegt wurde. Die React Native App kann einfach auf verschiedene Systeme portiert werden. Dabei müssen jedoch die jeweiligen Systemspezifischen Softwarespezifikationen beachtet werden. Dadurch können jedoch auch Features wie Benachrichtigungen oder Bluetooth-Verbindung ermöglicht werden.

% ---------------------------------------------------------------------------
\subsection{Pixel Streaming}
\label{subsec:Pixel Streaming}

Für \textbf{Pixel Streaming} ist eine Anbindungsstelle im Frontend
vorbereitet, die einen lokal verfügbaren Unreal-Engine-Stream einbetten kann.
Der entsprechende Backend-Endpunkt ist angelegt, jedoch wurden weder
Implementierung noch Tests durchgeführt. Dabei können komplexe Szenen einfach in einem Unreal Engine gebaut und getestet werden und im Anschluss in 
Diese und andere Erweiterungen werden in Abschnitt~\ref{chap:Diskussion}
weiter diskutiert.

\subsection{Raspberry Pi OS System}
\label{subsec:RaspberryPiOs}
Das Entwickelte System basiert auf einem RaspberryPi OS Build. Dabei wurde das zuletzt veröffentlichte Bullseye System genutzt. Dieses ist das letzte, welches die AV-Funktionalität nativ unterstützt. Sollte das System jemals neu aufgesetzt werden müssen, kann dies in der GitHub Repository gefunden werden. Die konfigurierten Zugangsdaten zu dem RaspberryPi Os sind:

User: admin und Passwort:admin

Durch das Anschließen des Systems an einen Monitor über die HDMI Schnittstelle können weiter Systemfunktionen konfiguriert werden und zum Beispiel die WIFI Verbindung zu einem neuen Netzwerk konfiguriert werden. In folgenden Iterationen könnte dies auch per Bluetooth funktionieren. 